% ============================================================================
% Appendix D: Hyperparameter Ablation Studies
% ============================================================================
% Simplified to focus on significant findings only.

\section{Hyperparameter Ablation Studies}
\label{app:ablations}

The main results (\S\ref{sec:metalearning_cost}) use constant
exploration ($\alpha{=}2.0$), mixing parameter ($\gamma{=}0.05$), and
learning rate ($\eta{=}\sqrt{\ln K/T}{\approx}0.030$).
This appendix validates these choices.

% ============================================================================
% D.1: Exploration Strategy (Alpha) Ablation
% ============================================================================

\subsection{Exploration Strategy}
\label{sec:alpha_ablation}

We evaluate four $\alpha$-configurations on the 750-prompt holdout set
($N{=}20$ seeds, Corralling with $\eta{=}0.030$, $\gamma{=}0.05$).

\begin{table}[h]
\centering
\caption{%
    \textbf{Exploration Strategy Ablation.}
    Two statistically distinguishable groups emerge:
    decay-based (${\sim}41$) and constant-based (${\sim}48$).
}
\label{tab:alpha_ablation}
\small
\begin{tabular}{@{}lcccc@{}}
\toprule
\textbf{Configuration} & \textbf{Warmup} & \textbf{Tabula Rasa} & \textbf{Regret} & \textbf{$p$ vs constant} \\
\midrule
\textbf{Homogeneous Decay} & decay & decay & \textbf{40.6 $\pm$ 4.4} & $<$0.001 \\
Mixed & constant & decay & 41.9 $\pm$ 4.4 & $<$0.001 \\
\midrule
Reversed Mixed & decay & constant & 46.8 $\pm$ 4.6 & ns \\
Homogeneous Constant & constant & constant & 48.2 $\pm$ 5.1 & --- \\
\bottomrule
\end{tabular}
\end{table}

\paragraph{Finding.}
Decay-based configurations (40.6--41.9) significantly outperform
constant exploration (46.8--48.2) on this static benchmark ($p{<}0.001$,
Wilcoxon signed-rank, Cohen's $d{\approx}1.6$).
Within each group, differences are not significant ($p{>}0.35$).

\paragraph{Mechanism.}
On a fixed dataset, decaying $\alpha$ allows experts to commit to
learned preferences once sufficient evidence accumulates.
The benefit is concentrated in the tabula rasa expert: decaying its
$\alpha$ captures most of the improvement (41.9 vs 48.2), because it
starts from zero and benefits most from transitioning to exploitation.

\paragraph{Important Caveat.}
This finding is specific to \emph{static} benchmarks.  In
non-stationary environments (\S\ref{sec:catastrophic}), premature
exploitation prevents detection of model degradation.
We therefore recommend constant $\alpha{=}2.0$ as the safe default,
accepting a 16\% regret premium as the cost of change detection.

\begin{figure}[h]
    \centering
    \includegraphics[width=0.9\textwidth]{experiments_v1/03_figure/results/ablation/figure_alpha_ablation.pdf}
    \caption{%
        \textbf{Exploration Strategy Ablation ($N{=}20$ seeds).}
        \textbf{(A)} Mean regret with $\pm$1~std.
        Decay configurations cluster around 41; constant around 48.
        \textbf{(B)} Per-seed scatter confirms tight clustering within
        each group and separation between groups.
    }
    \label{fig:alpha_ablation}
\end{figure}

% ============================================================================
% D.2: Mixing Parameter (Gamma) — Null Result
% ============================================================================

\subsection{Mixing Parameter}
\label{app:gamma_ablation}

The mixing parameter $\gamma$ guarantees minimum expert selection
probability $P(\text{expert}) \geq \gamma/K$, preventing expert death.
With the theoretically optimal $\eta{\approx}0.030$, $\gamma$ has
\textbf{negligible impact}: across
$\gamma \in \{0.00, 0.05, 0.10, 0.20\}$, cumulative regret ranges
from 47.9 to 48.6 (total span: 0.7 regret, well within one standard
deviation; all pairwise $p{>}0.5$, Wilcoxon signed-rank).

This null result is expected: at $\eta{\approx}0.030$, expert weights
evolve gradually from 50/50 to ${\sim}$71/29 over 750~rounds, never
approaching zero.  The $\gamma$-mixing floor is never activated.
We recommend $\gamma{=}0.05$ as a zero-cost safety default for longer
production deployments where weights might reach more extreme values.

% ============================================================================
% D.3: Prior-Trust Bias (Initial Weight) Sweep
% ============================================================================

\subsection{Prior-Trust Bias}
\label{app:initial_weight}

Section~\ref{sec:metalearning_cost} shows that Corralling incurs 60\%
overhead at $\alpha{=}0$ because the meta-learner starts with uniform
weights (50/50) and takes the full 750-round horizon to converge.
A natural question: can we reduce this overhead by \emph{biasing the
initial weights toward the warmup expert}?

We sweep the initial warmup weight $w_0 \in \{0.5, 0.6, 0.7, 0.8, 0.9\}$
across 7~corruption levels ($N{=}20$ seeds, $\eta{=}0.030$,
$\gamma{=}0.05$).  This encodes \emph{prior trust}: higher $w_0$
means the practitioner believes their priors are likely correct.

\begin{table}[h]
\centering
\caption{%
    \textbf{Prior-Trust Bias Trade-Off.}
    Higher initial warmup weight reduces regret when priors are correct
    ($\alpha{=}0$) but increases regret when priors are adversarial
    ($\alpha{=}1$).  $w_0{=}0.7$ provides a symmetric trade-off.
}
\label{tab:initial_weight}
\small
\begin{tabular}{@{}ccccc@{}}
\toprule
\textbf{$w_0$} & \textbf{Regret $\alpha{=}0$} & \textbf{Regret $\alpha{=}1$} &
\textbf{$\Delta$ $\alpha{=}0$} & \textbf{$p$ ($\alpha{=}1$)} \\
\midrule
0.5 (default) & 49.7 $\pm$ 5.1 & 59.9 $\pm$ 3.7 & --- & --- \\
\textbf{0.6} & \textbf{46.8 $\pm$ 4.8} & 60.5 $\pm$ 3.3 & $-$6\% & 0.49 (ns) \\
0.7 & 44.0 $\pm$ 4.7 & 65.6 $\pm$ 5.9 & $-$11\% & 0.014* \\
0.8 & 40.8 $\pm$ 3.8 & 75.6 $\pm$ 7.2 & $-$18\% & $<$0.001*** \\
0.9 & 36.6 $\pm$ 3.3 & 89.1 $\pm$ 7.0 & $-$26\% & $<$0.001*** \\
\bottomrule
\end{tabular}
\end{table}

\paragraph{Finding.}
$w_0{=}0.6$ is essentially a \textbf{free lunch}: 6\% lower regret at
$\alpha{=}0$ ($p{=}0.028$) with no significant cost at $\alpha{=}1$
($p{=}0.49$).  At $w_0{=}0.7$, the trade-off becomes approximately
symmetric: 11\% improvement at $\alpha{=}0$, 9\% cost at $\alpha{=}1$.
Beyond $w_0{=}0.8$, the cost grows superlinearly---the meta-learner
cannot recover fast enough from a biased starting point when the
warmup expert is actively harmful.

\paragraph{Mechanism.}
Setting $w_0{=}0.7$ gives the meta-learner a head start equivalent to
${\sim}28$ rounds of evidence that the warmup expert is better
(via the relationship $w_0/w_1 = e^{-\eta \Delta L}$).
When priors are correct, this skips most of the convergence cost.
When priors are wrong, the meta-learner must overcome this bias
\emph{plus} the original convergence, producing asymmetric degradation
at high bias values.

\paragraph{Recommendation.}
We recommend $w_0{=}0.7$ as the default for practitioners who have
\emph{any} prior information (the common case).
Use $w_0{=}0.5$ (uniform) only when prior quality is completely unknown
or when the maximum-safety guarantee is required.

\begin{figure}[h]
    \centering
    \includegraphics[width=0.95\textwidth]{experiments_v1/03_figure/results/initial_weight_sweep/figure_initial_weight_sweep.pdf}
    \caption{%
        \textbf{Prior-Trust Bias Trade-Off ($N{=}20$ seeds).}
        \textbf{(A)} Regret vs corruption at each initial warmup weight.
        Higher $w_0$ shifts Corralling's curve down at low corruption
        (closer to warmup-only) and up at high corruption (further from
        tabula rasa).  The curves fan out above $\alpha{=}0.5$.
        \textbf{(B)} Pareto frontier: each $w_0$ defines a point in
        (best-case regret, worst-case regret) space.
        Warmup-only (top-left) has the best best-case but worst
        worst-case; tabula rasa (bottom-right) is the opposite.
        Corralling traces a frontier between them, with $w_0$
        controlling the operating point.
    }
    \label{fig:initial_weight_sweep}
\end{figure}
